% ======================================================================
\section{Rapidity logarithms: BFKL and JIMWLK}\label{sec:rapidity}
% ======================================================================

We now evaluate the rows of Table~\ref{tab:rows}. In this section we collect all terms of the cross section at order $g^4$ that are enhanced by a logarithm of the longitudinal cut-offs, and show that they are the leading-logarithmic evolution of the leading-order (LO) spectrum: the BFKL evolution of the projectile charge correlator over the rapidity interval between the projectile and the measured gluon, and the JIMWLK evolution of the target operator over the interval between the measured gluon and the target. We first identify the regions of the longitudinal integrals that produce the logarithms and follow one row in detail. We then show how the terms with four and with three charges disappear from the logarithmic terms, for which the non-commutativity of the charges is essential, and identify the remaining two-charge coefficients.

\subsection{Origin of the logarithms}\label{sec:loglong}

Rapidity logarithms arise from longitudinal integrals of the form $\int dp^+/p^+$ over the momentum of the second gluon, real (integrated over in Groups~II and~III) or virtual (inside the coefficients $\mf N_2$, $\mf A^{(3)}$, $\mf B_3$ and in the interference terms of Group~I). In every row the integrand behaves as $1/p^+$ in two regions: for $k^+\ll p^+<\vee$, where the second gluon is faster than the measured one and belongs to the projectile, and for $\Lambda<p^+\ll k^+$, where it is slower and is resolved by the measured gluon as part of the target. The first region produces $\log(\vee/k^+)$, the second $\log(k^+/\Lambda)$. Terms that do not depend on $k^+$ at all, in particular the normalisation $\mf N_2$ and the products of four single-emission amplitudes, carry the full logarithm
\begin{align}\label{eq:logsplit}
\log\frac{\vee}{\Lambda}=\log\frac{\vee}{k^+}+\log\frac{k^+}{\Lambda}\,,
\end{align}
which we split at the longitudinal momentum $k^+$ of the produced gluon. With this split every logarithmic term is assigned either to the projectile or to the target side. Terms proportional to $\log(1-k^+/\vee)$, which arise from the upper limit $\vee-k^+$ of the integrals with a soft parent gluon, vanish for $\vee\to\infty$ and are dropped.

\paragraph{Which region produces which logarithm.} The assignment is summarised in Table~\ref{tab:logregions}. It is a statement about kinematics: a gluon that is faster than the measured one is formed earlier and is, for the measured gluon, part of the projectile; a gluon that is slower is formed later, is emitted coherently by the valence charges and by the measured gluon, and is indistinguishable from a gluon of the target. The coefficients decide how large the logarithm is, the limits of the $p^+$ integral decide which logarithm it is, and these limits follow from the kinematic restrictions of Section~\ref{sec:groups}: the upper limit is $\vee$ for an independently emitted gluon and $\vee-k^+$ where a soft parent of momentum $p^++k^+$ occurs, and the difference between the two is the power correction $\log(1-k^+/\vee)$.
\begin{table}[t]
\centering
\renewcommand{\arraystretch}{1.25}
\footnotesize
\begin{tabular}{@{}>{\raggedright\arraybackslash}p{0.22\textwidth}>{\raggedright\arraybackslash}p{0.33\textwidth}>{\raggedright\arraybackslash}p{0.14\textwidth}>{\raggedright\arraybackslash}p{0.22\textwidth}@{}}
\toprule
region of the second gluon & its role & logarithm & evolution\\
\midrule
$k^+\ll p^+<\vee$ & faster than the measured gluon: an additional source in the projectile & $\log(\vee/k^+)$ & projectile, BFKL (Section~\ref{sec:BFKL})\\
$\Lambda<p^+\ll k^+$ & slower: emitted coherently by the charges and by the measured gluon, scatters with them & $\log(k^+/\Lambda)$ & target, JIMWLK (Section~\ref{sec:JIMWLK})\\
no dependence on $k^+$ ($\mf N_2$, four single emissions) & both of the above & $\log(\vee/\Lambda)$, split by Eq.~\eqref{eq:logsplit} & both\\
$p^+\sim k^+$ & comparable momenta & none & DGLAP, RC and finite terms\\
\bottomrule
\end{tabular}
\caption{The regions of the longitudinal momentum $p^+$ of the second gluon and the terms they produce.}
\label{tab:logregions}
\end{table}

\paragraph{The lists.} For each of the two logarithms we collect the coefficients of all rows in a list, the BFKL list for $\log(\vee/k^+)$ and the JIMWLK list for $\log(k^+/\Lambda)$; the names anticipate the result of this section. A row with the cut-offs minus its entries in the two lists, multiplied by the logarithms, has a finite limit $\Lambda\to0$, $\vee\to\infty$, which is the finite term of the row (Section~\ref{sec:finite}). Section~\ref{sec:ex_III_V} shows how the entries arise for one row.

\paragraph{Order of the discussion.} By the order of the charges, the logarithmic terms are quartic, cubic or quadratic in $\rho$ (Table~\ref{tab:rows}). We discuss them in this order, because the reordering of the non-commuting charges feeds terms with more charges into terms with fewer.

\subsection{A worked row: Row III.V}\label{sec:ex_III_V}

Every row is evaluated by the same sequence of steps: the coefficients of Appendix~\ref{app:coefficients} are inserted into the product of amplitudes that defines the row, with the conjugated factor to the left and at the primed points; Wilson lines at the same point are combined with $U^TU=1$ and $f^{abc}U^{bd}U^{ce}=U^{ah}f^{hde}$, while the charges are not moved; the $\delta$-functions of $\mf C_1$ and the polarisation sums are done; the $p^+$ integral is written in terms of a momentum fraction, and its end points, which give the rapidity logarithms, are taken out; and the transverse integrals that can be done in closed form are done. We follow these steps for Row~III.V, $-\mathcal A^\dagger_{CA}\mathcal A^{(1)}_{AA}+{\rm c.c.}$, which contributes to both logarithms and to the finite terms. In the conjugate amplitude the two final-state gluons are the daughters of a LO gluon that has split after the collision; in the amplitude they are two independently emitted gluons. The row is cubic in the charges, with one charge at $\v x'$ from the conjugate amplitude and two charges at $\v x$ and $\v y$ from the amplitude, and the parent gluon of the splitting is a soft mode, $\Lambda<p^+<\vee-k^+$.

\paragraph{Coefficients and colour.} Inserting $\mf C_1$, Eq.~\eqref{eq:C1X}, and $\mf A$, Eq.~\eqref{eq:A1X}, into Eq.~\eqref{eq:GroupIIIexplicit},
\begin{align}\label{eq:ex3_start}
\frac{d^3N^{\rm III.V}}{d^3k}=\ &\frac{1}{(2\pi)^3}\,\frac{ig^4f^{c'd'a}}{4\pi^4}\,\frac{1}{k^{+2}}\int_\Lambda^{\vee-k^+}\frac{dp^+}{2\pi}\int_{w,w',y',z,y,x,x'}e^{-i\v k\cdot(\v w-\v w')}\notag\\
&\times\Kv^m(w'-z)\,\Kv^{k'}(x'-y')\,\Kv^i(y-w)\,\Kv^j(x-z)\;\,\delta^{(2)}\Big[\v w'-\frac{(p^++k^+)\v y'-p^+\v z}{k^+}\Big]\notag\\
&\times\Big[\delta_{k'm}\delta_{ij}-\frac{p^++k^+}{p^+}\,\delta_{jm}\delta_{ik'}-\frac{p^++k^+}{k^+}\,\delta_{im}\delta_{jk'}\Big]\notag\\
&\times\big[U^{c'e'}(y')-U^{c'e'}(x')\big]U^{d'b}(z)\,\rho^{e'}(x')\notag\\
&\times\big[U^{ad}(y)-U^{ad}(w)\big]U^{ce}(x)\,U^{cb}(z)\,\rho^e(x)\rho^d(y)+{\rm c.c.}
\end{align}
The single structure constant is that of $\mf C_1$. Since $U^{d'b}(z)U^{cb}(z)=\delta^{d'c}$, the integrated gluon does not see the target in the product, and the colour operator is
\begin{align}\label{eq:ex3_colour}
f^{mkl}\big[U^{ma}(y')-U^{ma}(x')\big]\,U^{kb}(x)\,\big[U^{lc}(y)-U^{lc}(w)\big]\;\rho^a(x')\,\rho^b(x)\,\rho^c(y)\,,
\end{align}
of the form $[U-U]\times U\times[U-U]\times f$ that is common to all cubic rows (Section~\ref{sec:fin_singlet}). The charges stand in the order conjugate amplitude, then amplitude, and within the amplitude in the order of $\mathcal A^{(1)}_{AA}$.

\paragraph{Recoil.} The $\delta$-function fixes $\v w'=\v y'+\frac{p^+}{k^+}(\v y'-\v z)$: the measured daughter is displaced from the parent, away from the other daughter. The integral over $\v w'$ therefore replaces the phase by
\begin{align}\label{eq:ex3_phase}
e^{-i\v k\cdot(\v w-\v w')}\to e^{-i\v k\cdot(\v w-\v y')}\,e^{-i\v k\cdot(\v z-\v y')\,\xi}\,,\qquad\xi=\frac{p^+}{k^+}\,,
\end{align}
and $\Kv^m(w'-z)$ by $\frac{1}{1+\xi}\Kv^m(y'-z)$. The second factor of Eq.~\eqref{eq:ex3_phase} is the \emph{recoil phase}. It occurs in all rows in which the measured gluon is part of a two-gluon state.

\paragraph{Longitudinal integral.} After this step the three tensor structures come with the factors
\begin{align}\label{eq:ex3_three}
\frac{1}{1+\xi}\,\delta_{k'm}\delta_{ij}-\frac1\xi\,\delta_{jm}\delta_{ik'}-\delta_{im}\delta_{jk'}\,.
\end{align}
under $\int d\xi$, with $\Lambda/k^+<\xi<\vee/k^+-1$. The term $1/\xi$ is split as $\frac1\xi=\frac{1}{\xi(1+\xi)}+\frac{1}{1+\xi}$. The first part converges at large $\xi$ and is written with a plus prescription, Eq.~\eqref{eq:plus_xi}; the subtracted term is the value of the whole integrand at $\xi=0$, where the recoil phase is one and $\v y'=\v w'$, times
\begin{align}
\int_{\Lambda/k^+}^{\vee/k^+-1}\frac{d\xi}{\xi(1+\xi)}=\log\frac{k^+}{\Lambda}+\log\Big(1-\frac{k^+}{\vee}\Big)+O\Big(\frac{\Lambda}{k^+}\Big)\,.
\end{align}
This is the entry of the row in the list of the terms proportional to $\log(k^+/\Lambda)$. The second part, $\frac{1}{1+\xi}$, is combined with the first term of Eq.~\eqref{eq:ex3_three}. The third term of Eq.~\eqref{eq:ex3_three} does not decrease with $\xi$. At large $\xi$ the recoil phase oscillates rapidly unless $\v y'\to\v z$, so that only this region contributes, and there the rest of the integrand can be evaluated at $\v y'=\v z$. We subtract and add the integrand with $\Kv(x'-y')\to\Kv(x'-z)$ and $U(y')\to U(z)$. In the added term the homogeneity of $\Kv$ gives
\begin{align}\label{eq:ex3_FT}
\int_{y'}e^{-i\v k\cdot(\v w-\v y')}\,e^{-i\v k\cdot(\v z-\v y')\,\xi}\,\Kv^m(y'-z)=\frac{1}{1+\xi}\int_{w'}e^{-i\v k\cdot(\v w-\v w')}\,\Kv^m(w'-z)\,,
\end{align}
so that it is $-\int d\xi/(1+\xi)=-\log(\vee/k^+)+O(\Lambda/k^+)$ times a structure in which the measured gluon in the conjugate amplitude is emitted from the position $\v z$ of the second gluon. This is the entry of the row in the list of the terms proportional to $\log(\vee/k^+)$: the second gluon is a fast gluon of the projectile that acts as a source for the measured one. The two terms with $\frac{1}{1+\xi}$ also produce logarithms of $\vee$, which cancel each other. With the plain split $[\frac1\xi]_+$ instead of $[\frac{1}{\xi(1+\xi)}]_+$ the remainder would still grow logarithmically with $\vee$; the split used here takes out the complete $\log(\vee/k^+)$ of the row.

\paragraph{The remainder.} What is left for $\Lambda\to0$ and $\vee\to\infty$ is the finite term of the row: the combination of the terms with $\frac{1}{1+\xi}$, the plus-prescription term, and the third term of Eq.~\eqref{eq:ex3_three} with its subtraction, which depends on $\xi$ only through the phase and is defined as a distribution, Eq.~\eqref{eq:inst_xi}. These are the kernels~\eqref{eq:K-y_rm_III} and~\eqref{eq:K-z_rm_III} of Section~\ref{sec:fin_rows}, and the finite term is given in Eqs.~\eqref{eq:S-III.V-3} and~\eqref{eq:S-III.V-2}. A single interference row thus contributes to both evolutions and to the finite terms, and the split between them is made by the behaviour in $\xi$ of three tensor structures.

\subsection{Terms quartic in \texorpdfstring{$\rho$}{rho}}\label{sec:fourrho}

Terms with four charges arise in Rows I.I, I.II (through the three-$\rho$ part $\mf A\mf N_2$ of $\mf A^{(3)}$), I.IV, II.I, II.II, II.III, III.I, III.II and III.IV, from the products of four single-emission coefficients $\mf A$, of $\mf A$ with $\mf N_2$, and of $\mf A$ with the $\{\rho,\rho\}$ part of $\mf B_2$. In all of them the longitudinal integral runs over the full interval $\Lambda<p^+<\vee$ and gives $\log(\vee/\Lambda)$, split according to Eq.~\eqref{eq:logsplit}.

\paragraph{Commuting charges.} The totally symmetrised four-$\rho$ part of their sum vanishes identically, i.e.\ the four-$\rho$ terms cancel for commuting charges. This cancellation has a simple origin. For commuting charges the part of $\Omega$ that is linear in $\rho$ exponentiates: to the order considered, $\Omega$ acting on the vacuum is the coherent state
\begin{align}\label{eq:coherent}
\Omega_{\rm coh}\ket0=\exp\Big\{\int\frac{d^3p}{(2\pi)^3}\Big[\mf A^b_j(p)\,a^{\dagger b}_j(p)-\mf A^{b\dagger}_j(p)\,a^b_j(p)\Big]\Big\}\ket0\,,
\end{align}
whose expansion to second order produces $\mf N_2$ and the $\{\rho,\rho\}$ part of $\mf B_2$, $\mf B_2\supset\frac12\mf A\mf A$. Since $\bar\Omega_{\rm coh}=S^\dagger\Omega_{\rm coh}S$ is again a displacement operator, $\bar\Omega^\dagger_{\rm coh}\Omega_{\rm coh}$ is, up to a phase, a displacement operator, and the expectation value of the number operator in the corresponding coherent state is the squared modulus of the displacement. The displacement is the LO amplitude~\eqref{eq:ALO}. Hence for commuting charges the spectrum is exactly the LO result, with no correction of order $g^4\rho^4$. Classical, commuting sources radiate a classical field, and to this order in the field of the projectile its amplitude is the LO one.

\paragraph{Two conditions.} Two features of the coefficients are necessary for this. First, the $\{\rho,\rho\}$ part of $\mf B_2$ must have support on the full square $\Lambda<k^+,p^+<\vee$, as in Eq.~\eqref{eq:B2X}; with the support $k^++p^+<\vee$ the four-$\rho$ terms of the different rows would carry different logarithms and would not cancel. Second, the three-$\rho$ part of $\mf A^{(3)}$ must be ordered as $\mf A\mf N_2$ (Section~\ref{sec:coefficients}), so that the normalisation terms combine into Eq.~\eqref{eq:Nterms}. We have verified the cancellation by an explicit term-by-term evaluation, with the charges represented by $SU(2)$ and $SU(3)$ matrices and the order of all operators kept.

\paragraph{Non-commuting charges.} For non-commuting charges the four-$\rho$ terms do not vanish as operators: rewriting each product in terms of totally symmetrised products leaves remainders with commutators. For two and three operators
\begin{align}\label{eq:symmetrise}
\rho_1\rho_2&=\tfrac12\{\rho_1,\rho_2\}+\tfrac12[\rho_1,\rho_2]\,,\notag\\
\rho_1\rho_2\rho_3&=S(\rho_1\rho_2\rho_3)+\tfrac12\Big(S\big([\rho_1,\rho_2]\rho_3\big)+S\big([\rho_1,\rho_3]\rho_2\big)+S\big([\rho_2,\rho_3]\rho_1\big)\Big)\notag\\&\quad+\tfrac16\Big(\big[[\rho_1,\rho_2],\rho_3\big]-\big[[\rho_2,\rho_3],\rho_1\big]\Big)\,,
\end{align}
where $S$ denotes the totally symmetrised product, $S(XY)=\frac12\{X,Y\}$, and similarly for four operators. By Eq.~\eqref{eq:rhoalgebra} each commutator reduces the number of charges by one. The four-$\rho$ terms therefore leave a three-$\rho$ remainder, multiplied by $\log(\vee/k^+)$ and by $\log(k^+/\Lambda)$, and a two-$\rho$ remainder, which is part of the BFKL and JIMWLK terms below. Since every four-$\rho$ term carries the full logarithm $\log(\vee/\Lambda)$ and nothing else, so do all its remainders: the four-$\rho$ rows have no finite part. This is why Rows~I.I, II.II, II.III and~III.IV do not appear in Section~\ref{sec:finite}. The situation can be summarised as a cascade,
\begin{align}\label{eq:cascade}
\underbrace{\rho\rho\rho\rho}_{\text{sym.\ part cancels}}\ \xrightarrow{\ [\rho,\rho]\ }\ \underbrace{\rho\rho\rho}_{\text{sym.\ part of the log terms cancels}}\ \xrightarrow{\ [\rho,\rho]\ }\ \underbrace{\rho\rho}_{\text{BFKL, JIMWLK}}\ \xrightarrow{\ [\rho,\rho]\ }\ \underbrace{\rho}_{\langle\rho\rangle=0}\,,
\end{align}
in which each arrow is one commutator, i.e.\ one factor $if^{abc}\delta^{(2)}$. The evolution of a two-charge correlator thus receives contributions from terms that, counted by the number of charges before reordering, have three and four charges; with commuting charges these contributions would be missing, and the coefficients of the logarithms would not be the BFKL and JIMWLK ones.

\subsection{Terms cubic in \texorpdfstring{$\rho$}{rho}}\label{sec:threerho}

Logarithmic terms with three charges arise directly in Rows I.II--I.V, II.I, III.I--III.III, III.V and III.VI, from the parts of $\mf B_2$, $\mf B_3$ and $\mf A^{(3)}$ that are linear or quadratic in $\rho$ multiplied by single emissions, and in addition from the reordering of the four-$\rho$ terms. They are of the form $if^{abc}\times$(real kernel)$\times$(Wilson lines)$\times\rho\rho\rho$, they change sign together with the sign of the exponent $e^{-i\v k\cdot(\v w-\v w')}$, and they are of the type of the ``first saturation corrections'' in the projectile.

\paragraph{The cancellation.} For each of the two logarithms, the totally symmetrised three-$\rho$ part of the sum of all these terms vanishes identically, for an arbitrary function of $\v w$ and $\v w'$ in place of the phase. We have organised the sum by transverse kernel and by the number of Wilson lines (two, three and four) and verified the cancellation in each class separately, analytically, at the operator level with explicit $SU(2)$ matrices for the charges on a transverse lattice, so that contact terms are included, and for the totally symmetrised colour tensors with random $SU(3)$ and $SU(4)$ Wilson lines (to $2\times10^{-15}$). Since the colour tensor itself vanishes, the cancellation holds for any symmetrised three-point function of the projectile, in particular for the $d^{abc}$ (odderon) one of Eq.~\eqref{eq:A2d}, whose projection we checked separately: no rapidity logarithm multiplies the three-point function. 

\paragraph{Charge conjugation.} That no logarithm can multiply $\sigma$ follows from charge conjugation. C acts as $\rho^at^a\to-(\rho^at^a)^T$, i.e.\ $\rho^a\to-\eta_a\rho^a$ with $(t^a)^T=\eta_at^a$, on the soft gluons as $A^a\to-\eta_aA^a$ and on the target as $V\to V^*$, $U^{ab}\to\eta_a\eta_bU^{ab}$; since $\eta_a\eta_b\eta_cf^{abc}=-f^{abc}$ and $\eta_a\eta_b\eta_cd^{abc}=d^{abc}$, it preserves Eq.~\eqref{eq:rhoalgebra} and all couplings, and the spectrum is invariant under the joint transformation. Under C, $n$ is even and $\sigma$ odd, and the Wilson-line traces with $T^e$ are even and those with $D^h$ odd (Section~\ref{sec:fin_singlet}). By the factorisation of Section~\ref{sec:proofsketch}, the coefficient of $\log(\vee/k^+)$ is the LO spectrum with the projectile correlator replaced by its one-step evolution, which does not involve the target: a three-charge term can enter it only through a $c$-number invariant tensor, $d^{abc}$ or $f^{abc}$, and C-invariance of the evolution admits only $f^{abc}$, whose contraction with Eq.~\eqref{eq:A2d} vanishes; the target enters through the C-even $\mathcal O_{\rm LO}$. The coefficient of $\log(k^+/\Lambda)$ is $-H_{\rm JIMWLK}$ acting on the target operator of the LO spectrum, in which the projectile enters only through $n$. Hence neither logarithm can multiply $\sigma$; a dependence on $\sigma$ requires a C-odd target structure and is non-logarithmic (Section~\ref{sec:fin_rows}). 

Charge conjugation alone does not forbid terms $\sigma\times$(C-odd target structure), which are exactly the symmetrised three-$\rho$ finite terms; the factorised form of the logarithms is needed in addition. The explicit cancellation confirms the argument and is stronger: individual terms have $d^{abc}$ projections of order one, with C-odd target factors, and they cancel in the sum.

\paragraph{Remainders.} As operators the three-$\rho$ terms leave, by Eq.~\eqref{eq:symmetrise}, two-$\rho$ and one-$\rho$ remainders. The two-$\rho$ remainders are part of the BFKL and JIMWLK terms. The one-$\rho$ remainders are contact terms $\propto f^{abc}\rho^c(\v x)\delta^{(2)}(\v x-\v y)$, which do not contribute in the colour-singlet form.

The non-logarithmic three-$\rho$ terms do not cancel. They belong to the finite terms (Section~\ref{sec:finite}).

\subsection{The BFKL term}\label{sec:BFKL}

Collecting all two-$\rho$ terms proportional to $\log(\vee/k^+)$, the genuine ones of the rows of Table~\ref{tab:rows} (I.II--I.IV, II.I, II.IV, III.I--III.III, III.V, III.VI) and the remainders of the reordered three- and four-$\rho$ terms, we obtain the BFKL term. For ordered charges of an arbitrary projectile it is given in Eq.~\eqref{eq:BFKLrho} of Appendix~\ref{app:finite}. For a colour-singlet projectile, Eq.~\eqref{eq:singlet}, it becomes
\begin{align}\label{eq:BFKLsinglet}
\log\frac{\vee}{k^+}\,\frac{d^{3}N_{\rm BFKL}}{d^{3}k} & = \frac{1}{(2\pi)^{3}}\,\frac{g^{4}}{16\pi^{5}}\,\frac{1}{k^{+}}\,\frac{1}{N^{2}_{c}-1}\,\log\Big[\frac{\vee}{k^{+}}\Big]\,\int_{x,y,z,w,w'}e^{-ik\cdot (w-w')}\notag\\&\times\langle\rho^{f}(x)\rho^{f}(y) \rangle\,\Bigg[ \notag\\ & \mathcal K(x-w)\cdot\mathcal K(y-w')\,K(x,y,z)\Big[2 N_{c}\mathcal O_{\rm LO}(x,y;w,w')\Big]\notag\\
&+\mathcal K(x-w)\cdot\mathcal K(y-w')\,K(x,x,z)\Big[-N_{c}\mathcal O_{\rm LO}(x,y;w,w')\Big]\notag\\
&+\mathcal K(x-w)\cdot\mathcal K(y-w')\,K(y,y,z)\Big[-N_{c}\mathcal O_{\rm LO}(x,y;w,w')\Big]\notag\\
&+K(w,w',z)\,K(x,y,z)\Big[2 N_{c}\mathcal O_{\rm LO}(z,z;w,w')\Big]\notag\\
&+\mathcal K(x-w')\cdot\mathcal K(w-z)\,K(x,y,z)\Big[2 N_{c}\mathcal O_{\rm LO}(z,x;w,w')\Big]\notag\\
&+\mathcal K(x-w')\cdot\mathcal K(w-z)\,K(y,y,z)\Big[-N_{c}\mathcal O_{\rm LO}(z,x;w,w')\Big]\notag\\
&+\mathcal K(x-w)\cdot\mathcal K(w'-z)\,K(x,y,z)\Big[2 N_{c}\mathcal O_{\rm LO}(x,z;w,w')\Big]\notag\\
&+\mathcal K(x-w)\cdot\mathcal K(w'-z)\,K(y,y,z)\Big[-N_{c}\mathcal O_{\rm LO}(x,z;w,w')\Big]\Bigg]
\end{align}
with $K(\v x,\v y,\v z)$ of Eq.~\eqref{eq:kernels} and $\mathcal O_{\rm LO}$ of Eq.~\eqref{eq:OLO}. As for LO, the sign of the exponent is immaterial in this form.

Eq.~\eqref{eq:BFKLsinglet} is the LO spectrum~\eqref{eq:LOsinglet} with the charge correlator replaced by its one-step BFKL evolution over the rapidity interval $\log(\vee/k^+)$:
\begin{align}\label{eq:BFKLasLO}
\log\frac{\vee}{k^+}\,\frac{d^3N_{\rm BFKL}}{d^3k}=\log\frac{\vee}{k^+}\;\frac{d^3N_{\rm LO}}{d^3k}\bigg|_{n(\v x,\v y)\to\frac{d}{dY}n(\v x,\v y)}\,,
\end{align}
where the evolution of the projectile charge correlator in the dilute regime is~\cite{Kovner:2006wr,Kovner:2006ge}
\begin{align}\label{eq:BFKLrhorho}
\frac{d}{dY}\big\langle\rho^f(\v x)\rho^f(\v y)\big\rangle=-\frac{\alpha_sN_c}{2\pi^2}\int_z\Big[&K_D(\v x,\v y,\v z)\,\big\langle\rho^f(\v x)\rho^f(\v y)\big\rangle
\notag\\&-2\,\delta^{(2)}(\v x-\v y)\int_{z'}K(\v z,\v z',\v x)\,\big\langle\rho^f(\v z)\rho^f(\v z')\big\rangle
\notag\\&
+\big(2K(\v y,\v z,\v x)-K(\v z,\v z,\v x)\big)\big\langle\rho^f(\v y)\rho^f(\v z)\big\rangle
\notag\\&+\big(2K(\v x,\v z,\v y)-K(\v z,\v z,\v y)\big)\big\langle\rho^f(\v x)\rho^f(\v z)\big\rangle\Big]\,,
\end{align}
with the dipole kernel $K_D(\v x,\v y,\v z)=K(\v x,\v x,\v z)-2K(\v x,\v y,\v z)+K(\v y,\v y,\v z)$. The first term of Eq.~\eqref{eq:BFKLrhorho} gives the first three lines of Eq.~\eqref{eq:BFKLsinglet} (virtual terms, the produced gluon is emitted by the original charges); the second term, in which the charge correlator is concentrated at the position $\v z$ of the new gluon, gives the fourth line (the produced gluon is emitted by the new gluon); and the last two terms give the remaining four lines (the produced gluon is emitted by the new gluon in the amplitude and by an original charge in the conjugate amplitude, or vice versa). In terms of the colour-singlet correlator of the dipole picture, Eq.~\eqref{eq:BFKLrhorho} is the linearised BK (dipole BFKL) equation; the check for a single colour dipole is described in Section~\ref{sec:checks}.

\paragraph{Interpretation.} The term describes a projectile that contains, at the moment of the collision, one more gluon than the valence charges, with a longitudinal momentum between $k^+$ and $\vee$. Being faster than the measured gluon, this gluon is a source for it, like a valence charge. This is why the Weizs\"acker--Williams field appears with the argument $\v w-\v z$ in the last five lines of Eq.~\eqref{eq:BFKLsinglet} and why the target operator is evaluated with one of the source positions replaced by $\v z$. The first three lines are the corresponding virtual corrections; together they are $-N_cK_D(\v x,\v y,\v z)\,\mathcal O_{\rm LO}(\v x,\v y;\v w,\v w')$. The logarithm is $\log(\vee/k^+)$ and not $\log(\vee/\Lambda)$ because only gluons faster than the measured one can act as its sources.

\subsection{The JIMWLK term}\label{sec:JIMWLK}

The two-$\rho$ terms proportional to $\log(k^+/\Lambda)$, the genuine ones of Rows I.II, I.V, II.I, II.IV and III.III and the remainders of the reordered three- and four-$\rho$ terms, add up to the JIMWLK term, which is given for ordered charges in Eq.~\eqref{eq:JIMWLKrho} of Appendix~\ref{app:finite}. For a colour-singlet projectile it is
\begin{align}\label{eq:JIMWLKsinglet}
    \log\frac{k^+}{\Lambda}\,\frac{d^{3}N_{\rm JIMWLK}}{d^{3}k} & = \frac{1}{(2\pi)^{3}}\,\frac{g^{4}}{16\pi^{5}}\,\frac{1}{k^{+}}\,\frac{1}{N^{2}_{c}-1}\,\log{\left(\frac{k^{+}}{\Lambda}\right)}\,\int_{x,y,z,w,w'}e^{-ik\cdot(w-w')}\notag\\&\times\langle\rho^{f}(x)\rho^{f}(y) \rangle\,\Bigg[\notag \\ & \mathcal K(x-w)\cdot\mathcal K(y-w')\,K(w,w',z)\Big[4 N_{c}{\rm Tr}\{U^\dagger(w)U(w')\}\notag\\
&-2{\rm Tr}\{T^eU^\dagger(w)U(w')T^eU^\dagger(w')U(z)\}\notag\\
&\quad-2{\rm Tr}\{T^eU^\dagger(w)U(w')T^eU^\dagger(z)U(w)\}\Big]\notag\\
&+\mathcal K(x-w)\cdot\mathcal K(y-w')\,K(x,w',z)\Big[-4 N_{c}{\rm Tr}\{U^\dagger(x)U(w')\}\notag\\
&+2{\rm Tr}\{T^eU^\dagger(w')U(x)T^eU^\dagger(x)U(z)\}\notag\\
&\quad+2{\rm Tr}\{T^eU^\dagger(w')U(x)T^eU^\dagger(z)U(w')\}\Big]\notag\\
&+\mathcal K(x-w)\cdot\mathcal K(y-w')\,K(x,y,z)\Big[4 N_{c}{\rm Tr}\{U^\dagger(x)U(y)\}\notag\\
&-2{\rm Tr}\{T^eU^\dagger(x)U(y)T^eU^\dagger(y)U(z)\}\notag\\
&\quad-2{\rm Tr}\{T^eU^\dagger(x)U(y)T^eU^\dagger(z)U(x)\}\Big]\notag\\
&+\mathcal K(x-w)\cdot\mathcal K(y-w')\,K(y,w,z)\Big[-4 N_{c}{\rm Tr}\{U^\dagger(w)U(y)\}\notag\\
&+2{\rm Tr}\{T^eU^\dagger(w)U(y)T^eU^\dagger(y)U(z)\}\notag\\
&\quad+2{\rm Tr}\{T^eU^\dagger(w)U(y)T^eU^\dagger(z)U(w)\}\Big]\notag\\
&+\mathcal K(x-w)\cdot\mathcal K(y-w')\,K(w',w',z)\Big[\notag\\
&\quad-2 N_{c}{\rm Tr}\{U^\dagger(w)U(w')\}+2 N_{c}{\rm Tr}\{U^\dagger(x)U(w')\}\notag\\
&-{\rm Tr}\{T^eU^\dagger(w')U(x)T^eU^\dagger(w')U(z)\}\notag\\
&\quad-{\rm Tr}\{T^eU^\dagger(w')U(x)T^eU^\dagger(z)U(w')\}\notag\\
&+{\rm Tr}\{T^eU^\dagger(w)U(w')T^eU^\dagger(w')U(z)\}\notag\\
&\quad+{\rm Tr}\{T^eU^\dagger(w)U(w')T^eU^\dagger(z)U(w')\}\notag\\
&+{\rm Tr}\{T^eU^\dagger(w')U(x)\}\,{\rm Tr}\{T^eU^\dagger(w')U(z)\}\notag\\
&\quad+{\rm Tr}\{T^eU^\dagger(w)U(w')\}\,{\rm Tr}\{T^eU^\dagger(w')U(z)\}\Big]\notag\\
&+\mathcal K(x-w)\cdot\mathcal K(y-w')\,K(w,w,z)\Big[\notag\\
&\quad-2 N_{c}{\rm Tr}\{U^\dagger(w)U(w')\}+2 N_{c}{\rm Tr}\{U^\dagger(w)U(y)\}\notag\\
&+{\rm Tr}\{T^eU^\dagger(w)U(w')T^eU^\dagger(w)U(z)\}\notag\\
&\quad+{\rm Tr}\{T^eU^\dagger(w)U(w')T^eU^\dagger(z)U(w)\}\notag\\
&-{\rm Tr}\{T^eU^\dagger(w)U(y)T^eU^\dagger(w)U(z)\}\notag\\
&\quad-{\rm Tr}\{T^eU^\dagger(w)U(y)T^eU^\dagger(z)U(w)\}\notag\\
&-{\rm Tr}\{T^eU^\dagger(w)U(w')\}\,{\rm Tr}\{T^eU^\dagger(w)U(z)\}\notag\\
&\quad+{\rm Tr}\{T^eU^\dagger(w)U(y)\}\,{\rm Tr}\{T^eU^\dagger(w)U(z)\}\Big]\notag\\
&+\mathcal K(x-w)\cdot\mathcal K(y-w')\,K(x,x,z)\Big[\notag\\
&\quad2 N_{c}{\rm Tr}\{U^\dagger(x)U(w')\}-2 N_{c}{\rm Tr}\{U^\dagger(x)U(y)\}\notag\\
&-{\rm Tr}\{T^eU^\dagger(w')U(x)T^eU^\dagger(x)U(z)\}\notag\\
&\quad-{\rm Tr}\{T^eU^\dagger(w')U(x)T^eU^\dagger(z)U(x)\}\notag\\
&+{\rm Tr}\{T^eU^\dagger(x)U(y)T^eU^\dagger(x)U(z)\}\notag\\
&\quad+{\rm Tr}\{T^eU^\dagger(x)U(y)T^eU^\dagger(z)U(x)\}\notag\\
&-{\rm Tr}\{T^eU^\dagger(w')U(x)\}\,{\rm Tr}\{T^eU^\dagger(x)U(z)\}\notag\\
&\quad-{\rm Tr}\{T^eU^\dagger(x)U(y)\}\,{\rm Tr}\{T^eU^\dagger(x)U(z)\}\Big]\notag\\
&+\mathcal K(x-w)\cdot\mathcal K(y-w')\,K(y,y,z)\Big[\notag\\
&\quad2 N_{c}{\rm Tr}\{U^\dagger(w)U(y)\}-2 N_{c}{\rm Tr}\{U^\dagger(x)U(y)\}\notag\\
&-{\rm Tr}\{T^eU^\dagger(w)U(y)T^eU^\dagger(y)U(z)\}\notag\\
&\quad-{\rm Tr}\{T^eU^\dagger(w)U(y)T^eU^\dagger(z)U(y)\}\notag\\
&+{\rm Tr}\{T^eU^\dagger(x)U(y)T^eU^\dagger(y)U(z)\}\notag\\
&\quad+{\rm Tr}\{T^eU^\dagger(x)U(y)T^eU^\dagger(z)U(y)\}\notag\\
&-{\rm Tr}\{T^eU^\dagger(w)U(y)\}\,{\rm Tr}\{T^eU^\dagger(y)U(z)\}\notag\\
&\quad+{\rm Tr}\{T^eU^\dagger(x)U(y)\}\,{\rm Tr}\{T^eU^\dagger(y)U(z)\}\Big]\Bigg]
\end{align} 
with $(T^e)_{bc}=-if^{ebc}$; the traces with $T^e$ are the adjoint form of the contractions $ffUUU$ of Eq.~\eqref{eq:JIMWLKrho}.

Eq.~\eqref{eq:JIMWLKrho} is the JIMWLK evolution of the target operator of the LO spectrum over the rapidity interval $\log(k^+/\Lambda)$. With the JIMWLK Hamiltonian in the normalisation of~\cite{Lublinsky:2016meo},
\begin{align}\label{eq:HJIMWLK}
H_{\rm JIMWLK}=\frac{\alpha_s}{2\pi^2}\int_{x,y,z}K(\v x,\v y,\v z)\Big[J_L^e(\v x)J_L^e(\v y)+J_R^e(\v x)J_R^e(\v y)-2\,c^{ef}(\v z)\,J_L^e(\v x)J_R^f(\v y)\Big]\,,
\end{align}
where $J_L^e(\v x)$ and $J_R^e(\v x)$ act on a Wilson line at $\v x$ as $J_L^eU=T^eU$ and $J_R^eU=UT^e$, $c^{ef}(\v z)$ is defined by $U^\dagger(\v z)T^eU(\v z)=c^{ef}(\v z)T^f$, and $\frac{d}{dY}=-H_{\rm JIMWLK}$, we find
\begin{align}\label{eq:JIMWLKasLO}
\log\frac{k^+}{\Lambda}\,\frac{d^3N_{\rm JIMWLK}}{d^3k}=-\log\frac{k^+}{\Lambda}\;H_{\rm JIMWLK}\,\frac{d^3N_{\rm LO}}{d^3k}\,,
\end{align}
where $H_{\rm JIMWLK}$ acts on the Wilson lines of Eq.~\eqref{eq:LO} only. The prefactors match since $\frac{g^4}{16\pi^5}\big/\frac{g^2}{2\pi^2}=\frac{\alpha_s}{2\pi^2}$. Because $\mathcal O_{\rm LO}$ is bilinear in $U$, $H_{\rm JIMWLK}$ acting on it can be evaluated exactly, for an arbitrary (not necessarily colour-singlet or symmetric) ordered correlator $\langle\rho^{b'}(\v x')\rho^b(\v x)\rangle$; Eq.~\eqref{eq:JIMWLKrho} agrees with the result pointwise in all transverse coordinates, and Eq.~\eqref{eq:JIMWLKsinglet} with its colour-singlet projection (Section~\ref{sec:checks}). For a single adjoint dipole the action of $H_{\rm JIMWLK}$ reduces to the BK form, $-H_{\rm JIMWLK}{\rm Tr}[U^\dagger(\v a)U(\v b)]=\frac{\alpha_s}{\pi^2}\int_zK_D(\v a,\v b,\v z)\big[c^{ef}(\v z){\rm Tr}(U^\dagger(\v a)T^eU(\v b)T^f)-N_c{\rm Tr}(U^\dagger(\v a)U(\v b))\big]$.

\paragraph{Interpretation.} Here the additional gluon is slower than the measured one, $\Lambda<p^+<k^+$. It is emitted, before or after the collision, by any of the colour charges that are present at its rapidity: the valence charges at $\v x$, $\v x'$ and the measured gluon at $\v w$, $\v w'$. The LO operator~\eqref{eq:OLO} is a sum of four adjoint dipoles ${\rm Tr}\{U^\dagger(\v a)U(\v b)\}$, with $\v a\in\{\v x,\v w\}$ from the amplitude and $\v b\in\{\v y,\v w'\}$ from the conjugate amplitude, and each of them evolves with the dipole kernel $K_D(\v a,\v b,\v z)$; a third Wilson line $U(\v z)$ appears when the slow gluon crosses the target. The slow gluon is not measured and does not change the kinematics of the measured one; its only effect is to change the colour rotation of the four lines. The logarithm is $\log(k^+/\Lambda)$, the rapidity interval between the measured gluon and the slowest mode that is still counted as part of the projectile wave function.

\subsection{The logarithms as the evolution of the leading-order spectrum}\label{sec:proofsketch}\label{sec:logsummary}

The identifications~\eqref{eq:BFKLasLO} and~\eqref{eq:JIMWLKasLO} are the expansion to first order of the leading-logarithmic factorisation of the single inclusive spectrum~\cite{Kovner:2006wr,Kovner:2006ge,Gelis:2008rw}: the spectrum of a gluon at rapidity $y$ is the LO expression evaluated with the projectile evolved from its own rapidity $Y_P$ down to $y$ and with the target evolved from $y$ to its own rapidity,
\begin{align}\label{eq:LLfact}
\frac{d^3N}{d^3k}\bigg|_{\rm LL}=\Big\langle\,\frac{d^3N_{\rm LO}}{d^3k}\,\Big\rangle_{P,\,Y_P-y;\;T,\,y-Y_T}\,.
\end{align}
The structure of the calculation shows why this holds, and we sketch the argument. Consider first the region $p^+\gg k^+$. Here the second gluon is much faster than the measured one, and its emission and absorption amplitudes factorise from those of the measured gluon: in all coefficients the dependence on $k^+$ disappears at leading power in $k^+/p^+$, and the measured gluon is emitted by the total colour charge of the valence charges and of the fast gluon, $\rho^a(\v x)\to\rho^a(\v x)+$ (charge density of the gluon at $\v z$). The sum over the rows is then the LO spectrum evaluated in the projectile state that contains one more soft gluon, minus the probability that it is not there (the virtual terms with $\mf N_2$ and the self-energy parts of $\mf A^{(3)}$), i.e.\ the LO spectrum with the projectile correlator evolved by one step. For a dilute projectile this step is linear in $\langle\rho\rho\rangle$, Eq.~\eqref{eq:BFKLrhorho}; terms with three and four charges, which would describe the non-linear (KLWMIJ) evolution of a dense projectile, must and do cancel at this order in the symmetrised sense (Sections~\ref{sec:fourrho},~\ref{sec:threerho}). In the region $p^+\ll k^+$ the second gluon is much slower than the measured one. It is emitted coherently by the valence charges and by the measured gluon, and at leading power its emission does not change the momentum of the measured gluon. The rows then combine into the change of the scattering amplitudes of the valence charges and of the measured gluon due to one soft gluon that scatters together with them, i.e.\ into one step of the JIMWLK evolution of the Wilson lines of the LO amplitude and of its conjugate, including the virtual terms. Both statements are verified by the explicit results~\eqref{eq:BFKLrho}--\eqref{eq:JIMWLKsinglet}.

In summary, the rapidity-enhanced part of the $O(g^4)$ cross section is
\begin{align}\label{eq:logsummary}
\frac{d^3N}{d^3k}\bigg|_{\log}=\log\frac{\vee}{k^+}\,\frac{d^3N_{\rm BFKL}}{d^3k}+\log\frac{k^+}{\Lambda}\,\frac{d^3N_{\rm JIMWLK}}{d^3k}\,,
\end{align}
with the BFKL term~\eqref{eq:BFKLrho} and the JIMWLK term~\eqref{eq:JIMWLKrho}, and no term with three or four charges multiplies a rapidity logarithm. Individual rows contain also double logarithms, $\log^2(k^+/\Lambda)$, from the region in which the second gluon is both soft and far from the measured one in the transverse plane. They come from the recoil of the measured gluon in the rows with $\mf B_3$ and $\mf C_1$ and cancel in the sum over rows (Appendix~\ref{sec:fin_soft}).

\paragraph{Use of the result.} Eq.~\eqref{eq:logsummary} is the first term of the expansion of the leading-logarithmic formula~\eqref{eq:LLfact}. In an application the two logarithms are not kept as they stand: $\log(\vee/k^+)$ is absorbed into the evolution of the projectile correlator from the rapidity of the projectile down to that of the measured gluon, and $\log(k^+/\Lambda)$ into the evolution of the target average of the Wilson lines from its own rapidity up to that of the measured gluon; the LO term is then evaluated with the evolved averages, and the remaining terms of Eq.~\eqref{eq:result} are the genuine corrections of order $\alpha_s$. The cut-offs $\vee$ and $\Lambda$ thereby become the rapidity-factorisation scales of the projectile and of the target.

